\section{Numerical implementation and study design}
\label{sec:numerics}
\subsection{Geometry and carrier preparation}
A structured axisymmetric 5$^\circ$ wedge represents the 1.25 m column
(\cref{fig:geometry}). The steel and silica radii are 2.4 and 2.5 mm,
respectively. The baseline mesh has $384\times48\times1=18{,}432$
cells. Coarse and fine meshes have $256\times32\times1=8{,}192$ and
$576\times72\times1=41{,}472$ cells, refining axial and radial spacing
by a factor of 1.5 between successive levels. The source spans
$x=0.09$--0.105 m and the radial interval $r=0.25R$--$0.75R$.
Its position and selection are computational assumptions; the physical
boat is omitted from the carrier geometry. Full-column source and inlet
rates are multiplied by $5/360=1/72$ for the wedge.

\begin{figure}[tb]
\centering
\includegraphics[width=.94\textwidth]{experiment_schematic.pdf}
\caption{Computational geometry and prescribed source region. The source
is distributed in selected cells; the sketch does not represent a
resolved boat surface. Tube length and material-specific diameters are
retained, and the 5$^\circ$ wedge supplies the axisymmetric reduction.}
\label{fig:geometry}
\end{figure}

The inlet flow rates are converted to helium mass flow assuming a
controller reference of 298.15 K and 101,325 Pa. Full-column rates are
$2.727\times10^{-7}\units{kg\,s^{-1}}$ at 100 mL/min and
$1.227\times10^{-7}\units{kg\,s^{-1}}$ at 45 mL/min. The reference
conditions are not specified by the supplied paper. A separate study
uses 273.15 K at the same reference pressure. Helium uses the perfect-gas
law, constant heat capacity and Sutherland viscosity
(\cref{app:properties}). Steady laminar carriers are prepared with
\texttt{rhoSimpleFoam}; the reconstructed temperature profile remains
prescribed.

The converged face mass flux is retained instead of reconstructing it
from interpolated velocity. At these low mass flows, residual discrete
divergence is removed by a zero-initialised potential correction before
species transport. Carrier preparation rejects an original cell or global
flux defect exceeding $10^{-6}$ of inlet flow, or a flux/velocity
correction exceeding $10^{-5}$ of the corresponding scale. After
correction, inlet, global and cell continuity and impermeable-boundary
leakage are audited against $10^{-12}$ of inlet flow. This procedure
makes the frozen transport field discretely conservative; it is not a
new physical flow model or evidence of convergence of the velocity field.

\subsection{Discretisation and coupling}
The transient species equations use implicit Euler time integration,
first-order upwind advection and corrected Gauss linear diffusion. All
nine transported gases use the same advection scheme. The baseline time
step is 2 ms. Wall exchange is assembled implicitly with nonnegative
source and sink coefficients and shared inventory bounds. Configured
species correctors and active-bound guards reconcile the final gas
solution with the realised wall transfer. Gas re-speciation then
redistributes species under the local element constraints. This sequence
is first order in time; a conservative equilibrium projection alone
does not eliminate splitting or transport error.

The baseline configuration requests three species correctors with a
$10^{-16}$ tolerance and up to eight bound guards with a $10^{-13}$
guard tolerance. These are stopping settings, not guarantees that every
outer wall-regime iteration attains the requested tolerance. The final
transport defect, bound status and material ledger must be examined
independently (\cref{sec:conservation}). Additional short corrector
replays isolate late-time iteration sensitivity but do not establish
full-run iteration independence.

The baseline diffusivity is $D_i=10^{-4}\units{m^2\,s^{-1}}$ for every
gas. Three transport alternatives halve or double this value for all
gases, or replace only the PbI$_2$ coefficient by a temperature- and
pressure-dependent PbI$_2$--He estimate. The latter's collision
parameters are estimates, not independently measured binary transport
data. Accommodation variations change the two reversible iodide pairs
to $\alpha=0.1$ or 0.01, retaining the baseline metal-channel coefficients.

\subsection{Completed case matrix}
The 15 baseline cases comprise Q1 for each available configuration and
Q2 for the two LBE columns, at
$\beta\in\{0.001,0.01,0.1\}$ with two thermodynamic branches. The
second branch adds a constant 29 kJ/mol to BiI(g)'s standard Gibbs energy;
it is a sensitivity surrogate, not a second experimentally qualified
dataset. All completed calculations use the same solver executable,
identified by its SHA-256 in the saved reports.

\Cref{tab:study} describes the 64 follow-up cases. Mesh and time-step
checks cover pure PbI$_2$ and three steel Q2 cases: baseline
thermodynamics at 0.1\% and 1\% metal saturation and shifted BiI
thermodynamics at 1\%. Physical transport, accommodation, source-position
and flow-reference variations cover pure PbI$_2$, those two baseline
steel Q2 cases, and silica Q2 at 0.1\%. Extra source-loading cases use
both columns and both thermodynamic branches. Temperature-coordinate
variants use silica Q1 and silica Q2 at 0.1\%. This is a structured
set of conditional comparisons, not a factorial uncertainty study.
\begin{table}[tb]
\centering
\caption{Completed follow-up matrix. The baseline cases are additional
references. Numerical comparisons use one changed resolution at a time;
flow-stop cases continue a 60 s baseline to 80 s.}
\label{tab:study}
\small\begin{tabular}{lrl}
\toprule
Study & Cases & Variation \\
\midrule
Mesh and time step & 12 & Three grids; 2 and 1 ms \\
Gas diffusion & 12 & $D/2$, $2D$, PbI2--He correlation \\
Iodide accommodation & 8 & $\alpha=0.1$, 0.01 \\
Source position & 8 & $\Delta x_s=\pm1.5$ cm \\
Q2 metal loading & 12 & $\beta=0.002$, 0.003162, 0.005623 \\
Flow reference & 4 & 273.15 versus 298.15 K \\
Silica temperature coordinates & 4 & $\Delta x_T=-4.10$, $-4.65$ cm \\
Flow-stop relaxation & 4 & 60--80 s; fixed temperature \\
\bottomrule
\end{tabular}

\end{table}

The source-position alternatives move the release region by
$\pm1.5\units{cm}$ without adding a boat obstruction. Flow-stop cases
restart the saved 60 s fields, remove injection, set carrier flow to
zero and impose zero normal species flux at both ends. The original
nonuniform temperature profile stays fixed. Their comparisons therefore
measure redistribution of residual gas during a 20 s fixed-temperature
flow stop, not a clean-helium purge or a full cooling calculation.
